\section{Introduction}
\label{sec:introduction}

Spiking neural networks (SNNs) represent information through discrete events evolving in continuous time. Recent exact-gradient methods have established that this timing is amenable to gradient-based optimization. EventProp treats an SNN as a hybrid dynamical system and propagates adjoint variables through both continuous neuron dynamics and discontinuous spike events, yielding exact gradients almost everywhere without replacing the spike threshold by a surrogate derivative \cite{wunderlich2021eventprop}. DelGrad and delayed extensions of EventProp further show that transmission delays can be optimized directly and can improve temporal-task performance and parameter efficiency \cite{goeltz2025delgrad,meszaros2025delay}. These results motivate treating event time, rather than only synaptic weight, as a first-class computational degree of freedom.

The temporal response of the standard current-based leaky integrate-and-fire (LIF) neuron nevertheless remains strongly constrained by its dynamics. With an exponential synapse, the postsynaptic-potential kernel is an alpha-shaped double low-pass filtering response. This kernel provides a single rise-and-decay profile, and adjusting a weight changes the scale of this fixed temporal basis and indirectly moves a threshold crossing; it does not change the qualitative geometry of the basis itself. Consequently, an EventProp-trained LIF network can optimize spike timing while still operating within a restricted family of temporal responses.

Two lines of work suggest complementary ways to relax this restriction. First, complex-conjugate decay modes provide damped oscillatory memory. Such modes are central to structured state-space models for long-sequence processing \cite{gu2022s4}, and resonate-and-fire neurons use the same mathematical structure to obtain frequency selectivity and sensitivity to relative spike timing \cite{izhikevich2001resonate}. Recent trainable resonate-and-fire networks confirm that oscillatory state dynamics can improve long-range temporal processing and reduce the number of required spiking operations \cite{huber2025s5rf}. Damped resonant kernels have also been considered explicitly as synaptic basis functions in SNNs \cite{tapson2013skim}. Second, models with spike latency and networks with learnable transmission delays show that additional latency freedom can increase the set of temporal transformations available to an SNN \cite{susi2018lifl,goeltz2025delgrad,meszaros2025delay}. However learned delays are typically neuron static parameters. QIF \cite{wenig2026qif} and LIFL \cite{susi2018lifl} suggest more powerful neural state dependent delay, which [TODO later]. Meanwhile, existing approaches generally study oscillatory dynamics and propagation delays separately.

We propose the \emph{Temporal-enhanced Spiking} (TeS) neuron, an LIF-compatible model that combines a real-valued soma with two differentiable temporal mechanisms: an oscillatory synaptic state and a state-dependent transmission delay. The soma integrates the real component of a damped complex synaptic trace and otherwise retains the ordinary LIF threshold and reset. At each upward threshold crossing, the pre-reset voltage slope determines a bounded delay before the spike is delivered to downstream synapses. The complete forward model is given in Section~\ref{sec:method}.

The two mechanisms implement complementary temporal operations. The complex pole supplies a frequency-selective, phase-sensitive memory state, whereas the delay performs event-wise temporal warping conditioned on the source neuron's pre-spike dynamics. Their interaction is the central computational primitive of TeS: the delay determines the phase at which an incoming spike interacts with the postsynaptic oscillatory state. TeS therefore converts analog information that a conventional point neuron discards at threshold---the sharpness of the crossing---into the timing of an otherwise binary event.

The contribution of this work is the unified engineering realization of these mechanisms within an exact event-based learning framework. Specifically, we aim to establish:
\begin{enumerate}
    \item an EventProp-compatible hybrid formulation for a real LIF soma driven by damped oscillatory synapses;
    \item differentiable scheduling of future synaptic events through a per-spike delay conditioned on the source-neuron threshold slope;
    \item the corresponding adjoint flow, event-jump conditions, and parameter gradients, which are exact wherever threshold crossings are transverse and the event topology is unchanged; and
    \item an event-driven, neuromorphic-compatible realization using two real synaptic state variables, binary spike communication, and quantities locally available at spike emission.
\end{enumerate}
To our knowledge, TeS is the first continuous-time SNN framework to integrate damped oscillatory synaptic dynamics and per-spike state-dependent transmission delays while retaining EventProp-style exact gradients almost everywhere. We do not claim oscillatory synapses or state-dependent neural delays individually as new. Rather, the proposed novelty lies in their joint differentiable realization and in \emph{slope-conditioned phase routing}: the source state selects when a spike reaches a resonant postsynaptic state and therefore how that spike interferes with recent temporal context.

This construction is significant if the two mechanisms provide more than additive model capacity. We hypothesize that oscillatory filtering expands the temporal basis available to each neuron, state-dependent delay aligns events with task-relevant phases, and their interaction improves temporal expressiveness beyond either mechanism alone. The claim can be tested with parameter-matched ablations against ordinary LIF, oscillatory-synapse-only, state-dependent-delay-only, and static or freely trainable-delay networks. In addition to task performance, the evaluation should measure spike count, latency, gradient accuracy, robustness near shallow threshold crossings, delay quantization, analog parameter mismatch, and noise. This separates increased temporal expressiveness from gains caused only by additional parameters and directly tests whether TeS offers a useful primitive for event-driven neuromorphic systems.
